% Template for post or presentation abstracts submitted to ILAS 2022
% 1. Fill in the presenter's information (name, affiliation, email
% address) and talk information (title of presentation, abstract).
% 2. Compile to make sure there are no errors.
% 3. Save the .tex file for your abstract with a distinctive name,
% ideally "FamilynameGivenname.tex".
% 4. Upload your .tex file to https://clr.ie/129557
%       
% * Please keep your LaTeX commands simple, and avoid the use of
%    user-defined macros.
% * Include details of co-authors and funding acknowledgements after the abstract.
% * Upload only the LaTeX file (with .tex extension).
% * Questions: Contact Niall Madden (Niall.Madden@NUIGalway.ie)

\documentclass[a4paper]{amsart} 
\usepackage{amssymb} 
\usepackage{hyperref}
\pagestyle{empty}
\date{} 

%% IGNORE THE NEXT 4 LINES
\makeatletter % 
\def\labelenumi{\theenumi.}
\def\theenumi{\@arabic\c@enumi}
\makeatother
%% STOP IGNORING NOW

\begin{document}

\title{Gaussian Elimination in the Nineteenth Century}
\author{Joseph Grcar} %Presenting author only
\address{Castro Valley, California} % Affiliation of presenting author
\email{jfgrcar@gmail.com} % Email address of presenting author only

\maketitle

% The abstract  ***no more than one page***

Gauss, Doolittle and Cholesky used what we call Gaussian elimination to make socially valuable calculations in the nineteenth century.
Namely, they solved minimum norm problems for angle adjustments in surveying.

Their work is a good example of how numerical calculations evolve.
A social need arises for answers to a recurring, mathematical question, which may be formulated in various ways as  problems of a numerical character.
Gauss devised the mathematical formulation of angle adjustments, which did not change thereafter.

Because calculations were tediously done by hand in the nineteenth century, the manner of making a calculation depended strongly on the availability of computing aids. 
Gauss, Doolittle and Cholesky used as aids logarithms, multiplication tables and pinwheel calculators, respectively.

Legions of students learned to understand the mathematical formulation of adjustments and to be proficient in the calculations.
This curriculum was vocational training, at the university level, for an early form of data analytics.
For example, American philosopher Henry David Thoreau apparently learned to be a professional surveyor at his alma mater, Harvard College.

It is appropriate to trace the history of numerical mathematics in terms of the actual use.
Algorithms in the sense of abstract discoveries are not the principal focus.
Instead, algorithms are merely responses to varying problem formulations and varying aids to computation.
Some examples in the nineteenth and earlier centuries are predicting artillery trajectories and ocean tides, and in the twenty-first century, perhaps, constructing artificial intelligence.

The historical material in this talk is taken from \cite {Grcar-HM}, which is summarized in \cite {Grcar-Notices}.
See \cite {Goodman-BAM} for a note on Thoreau.

\begin{thebibliography}{1}
\bibitem{Grcar-HM}
Joseph Grcar.
\newblock How ordinary elimination became Gaussian elimination.
\newblock {\em  Historia Mathematica,}  38(2):163--218, 2011.

\bibitem{Grcar-Notices}
Joseph Grcar.
\newblock Mathematicians of Gaussian elimination.
\newblock {\em  Notices of the AMS,} 58(6):782--792, 2011.

\bibitem{Goodman-BAM}
Lawrence Goodman.
\newblock Thoreau's Math.
\newblock {\em Brown Alumni Magazine,} January/February 2007.
\end{thebibliography}


\end{document}


% Please leave the next bit alone  
%%% Local Variables: 
%%% mode: latex
%%% TeX-master: "../ILAS-2022"
%%% End:
